\documentclass[10pt,a4paper]{amsart}
\usepackage{amsmath,amssymb,amsthm,graphicx}
\usepackage{verbatim}
\newenvironment{reference}{\comment}{\endcomment}
%\newenvironment{reference}{}{}
\newenvironment{slogan}{\comment}{\endcomment}
\newenvironment{history}{\comment}{\endcomment}
% For commutative diagrams we use Xy-pic
\usepackage[all]{xy}

% We use 2cell for 2-commutative diagrams.
\xyoption{2cell}
% Theorem environments.
%
\theoremstyle{plain}
\newtheorem{theorem}[subsection]{Theorem}
\newtheorem{proposition}[subsection]{Proposition}
\newtheorem{lemma}[subsection]{Lemma}

\theoremstyle{definition}
\newtheorem{definition}[subsection]{Definition}
\newtheorem{example}[subsection]{Example}
\newtheorem{exercise}[subsection]{Exercise}
\newtheorem{situation}[subsection]{Situation}

\theoremstyle{remark}
\newtheorem{remark}[subsection]{Remark}
\newtheorem{remarks}[subsection]{Remarks}

\numberwithin{equation}{subsection}

% Macros
%
\def\lim{\mathop{\mathrm{lim}}\nolimits}
\def\colim{\mathop{\mathrm{colim}}\nolimits}
\def\Spec{\mathop{\mathrm{Spec}}}
\def\Hom{\mathop{\mathrm{Hom}}\nolimits}
\def\Ext{\mathop{\mathrm{Ext}}\nolimits}
\def\SheafHom{\mathop{\mathcal{H}\!\mathit{om}}\nolimits}
\def\SheafExt{\mathop{\mathcal{E}\!\mathit{xt}}\nolimits}
\def\Sch{\mathit{Sch}}
\def\Mor{\mathop{\mathrm{Mor}}\nolimits}
\def\Ob{\mathop{\mathrm{Ob}}\nolimits}
\def\Sh{\mathop{\mathit{Sh}}\nolimits}
\def\NL{\mathop{N\!L}\nolimits}
\def\CH{\mathop{\mathrm{CH}}\nolimits}
\def\proetale{{pro\text{-}\acute{e}tale}}
\def\etale{{\acute{e}tale}}
\def\QCoh{\mathit{QCoh}}
\def\Ker{\mathop{\mathrm{Ker}}}
\def\Im{\mathop{\mathrm{Im}}}
\def\Coker{\mathop{\mathrm{Coker}}}
\def\Coim{\mathop{\mathrm{Coim}}}

% Boxtimes
%
\DeclareMathSymbol{\boxtimes}{\mathbin}{AMSa}{"02}

%
% Macros for moduli stacks/spaces
%
\def\QCohstack{\mathcal{QC}\!\mathit{oh}}
\def\Cohstack{\mathcal{C}\!\mathit{oh}}
\def\Spacesstack{\mathcal{S}\!\mathit{paces}}
\def\Quotfunctor{\mathrm{Quot}}
\def\Hilbfunctor{\mathrm{Hilb}}
\def\Curvesstack{\mathcal{C}\!\mathit{urves}}
\def\Polarizedstack{\mathcal{P}\!\mathit{olarized}}
\def\Complexesstack{\mathcal{C}\!\mathit{omplexes}}
% \Pic is the operator that assigns to X its picard group, usage \Pic(X)
% \Picardstack_{X/B} denotes the Picard stack of X over B
% \Picardfunctor_{X/B} denotes the Picard functor of X over B
\def\Pic{\mathop{\mathrm{Pic}}\nolimits}
\def\Picardstack{\mathcal{P}\!\mathit{ic}}
\def\Picardfunctor{\mathrm{Pic}}
\def\Deformationcategory{\mathcal{D}\!\mathit{ef}}
\usepackage{cleveref}
\crefname{originalreference}{original reference}{original references}
\Crefname{originalreference}{Original reference}{Original references}
\makeatletter
\newlabel{criteria-lemma-flat-quotient-flat-presentation}{{criteria-lemma-flat-quotient-flat-presentation}{source}{}{sourcepointer-1}{}}
\newlabel{criteria-lemma-flat-quotient-flat-presentation@cref}{{[originalreference][0][]criteria-lemma-flat-quotient-flat-presentation}{[1][0][]source}{}{}{}}
\makeatother
\crefname{originalreference}{original reference}{original references}
\Crefname{originalreference}{Original reference}{Original references}
\makeatletter
\newlabel{spaces-groupoids-lemma-quotient-stack-2-cartesian}{{spaces-groupoids-lemma-quotient-stack-2-cartesian}{source}{}{sourcepointer-2}{}}
\newlabel{spaces-groupoids-lemma-quotient-stack-2-cartesian@cref}{{[originalreference][0][]spaces-groupoids-lemma-quotient-stack-2-cartesian}{[1][0][]source}{}{}{}}
\makeatother
\crefname{originalreference}{original reference}{original references}
\Crefname{originalreference}{Original reference}{Original references}
\makeatletter
\newlabel{spaces-morphisms-lemma-universally-closed-permanence}{{spaces-morphisms-lemma-universally-closed-permanence}{source}{}{sourcepointer-3}{}}
\newlabel{spaces-morphisms-lemma-universally-closed-permanence@cref}{{[originalreference][0][]spaces-morphisms-lemma-universally-closed-permanence}{[1][0][]source}{}{}{}}
\makeatother
\crefname{originalreference}{original reference}{original references}
\Crefname{originalreference}{Original reference}{Original references}
\makeatletter
\newlabel{stacks-properties-lemma-check-property-covering}{{stacks-properties-lemma-check-property-covering}{source}{}{sourcepointer-4}{}}
\newlabel{stacks-properties-lemma-check-property-covering@cref}{{[originalreference][0][]stacks-properties-lemma-check-property-covering}{[1][0][]source}{}{}{}}
\makeatother
\crefname{originalreference}{original reference}{original references}
\Crefname{originalreference}{Original reference}{Original references}
\makeatletter
\newlabel{stacks-properties-section-properties-morphisms}{{stacks-properties-section-properties-morphisms}{source}{}{sourcepointer-5}{}}
\newlabel{stacks-properties-section-properties-morphisms@cref}{{[originalreference][0][]stacks-properties-section-properties-morphisms}{[1][0][]source}{}{}{}}
\makeatother
\title{Properties of the diagonal read off a presentation of the stack}
\author{The Stacks Project Authors}
\date{stacks-project:stacks-morphisms.tex:local-snapshot-20261012; GNU FDL 1.2 or later}
\begin{document}
\maketitle
\noindent\textit{Editorial note.} All delivered mathematical text is the author's own, copied from the pinned Stacks Project chapter source sources/stacks-algebraic-stacks-morphisms/pinned-original-source.tex; no statement or proof step was rewritten, reordered or removed.\par\medskip
\noindent\textit{Editorial note.} Preamble provenance: the Stacks Project keeps one single shared preamble for every chapter; this card takes it directly from the staged sources/stacks-algebraic-stacks-morphisms/preamble.tex as three verbatim spans (lines 8-12, 14-18 and 154-217), so no derived file was written and nothing was added or redefined. The wrapper class and the base packages are supplied by the assembler.\par\medskip
\noindent\textit{Editorial note.} Bibliography provenance: neither the counted statement nor its proof contains a citation (this lane rescans the whole span before writing the card), so the card carries no bibliography and no entry was reconstructed.\par\medskip
\noindent\textit{Editorial note.} Selection provenance: this card was opened under the file-based control contract (CONTROL.md), with the mechanical pre-card gate run on its statement and proof interval before the card was written.\par\medskip
\noindent\textit{Allowed standard prerequisites (editorial).}\begin{itemize}
\item The description of the diagonal through a two-fibre-product square over a presentation. Used at the start of the counted proof (source lines 4094-4105). Expresses the diagonal of a presentation morphism through the diagonal of the stack.
\item The criterion for checking a property of a representable morphism along a smooth surjective cover. Invoked inside the counted proof (source lines 4106-4115) as a cross-chapter pointer. Transfers properties of the diagonal between the stack and its presentation.
\item The corresponding statements for properties of morphisms of algebraic spaces. Invoked at the end of the counted proof (source lines 4116-4122) as cross-chapter pointers. Supplies the equivalences at the level of the presentation.
\end{itemize}\medskip
\ifdefined\hypertarget\hypertarget{sourcepointer-1}{}\fi
\noindent\textit{Source reference.} Original reference criteria-lemma-flat-quotient-flat-presentation: sources/stacks-algebraic-stacks-morphisms/pinned-original-source.tex, source lines 3877--3877. Navigation pointer only: the counted text refers to the Stacks item criteria-lemma-flat-quotient-flat-presentation, which belongs to another chapter of the project and is not part of this card. The pointer carries the target's own label as its number because this chapter's staged snapshot does not contain the other chapter's tag table; its content is not reproduced here.\par\medskip
\ifdefined\hypertarget\hypertarget{sourcepointer-2}{}\fi
\noindent\textit{Source reference.} Original reference spaces-groupoids-lemma-quotient-stack-2-cartesian: sources/stacks-algebraic-stacks-morphisms/pinned-original-source.tex, source lines 4107--4107. Navigation pointer only: the counted text refers to the Stacks item spaces-groupoids-lemma-quotient-stack-2-cartesian, which belongs to another chapter of the project and is not part of this card. The pointer carries the target's own label as its number because this chapter's staged snapshot does not contain the other chapter's tag table; its content is not reproduced here.\par\medskip
\ifdefined\hypertarget\hypertarget{sourcepointer-3}{}\fi
\noindent\textit{Source reference.} Original reference spaces-morphisms-lemma-universally-closed-permanence: sources/stacks-algebraic-stacks-morphisms/pinned-original-source.tex, source lines 1056--1056. Navigation pointer only: the counted text refers to the Stacks item spaces-morphisms-lemma-universally-closed-permanence, which belongs to another chapter of the project and is not part of this card. The pointer carries the target's own label as its number because this chapter's staged snapshot does not contain the other chapter's tag table; its content is not reproduced here.\par\medskip
\ifdefined\hypertarget\hypertarget{sourcepointer-4}{}\fi
\noindent\textit{Source reference.} Original reference stacks-properties-lemma-check-property-covering: sources/stacks-algebraic-stacks-morphisms/pinned-original-source.tex, source lines 479--479. Navigation pointer only: the counted text refers to the Stacks item stacks-properties-lemma-check-property-covering, which belongs to another chapter of the project and is not part of this card. The pointer carries the target's own label as its number because this chapter's staged snapshot does not contain the other chapter's tag table; its content is not reproduced here.\par\medskip
\ifdefined\hypertarget\hypertarget{sourcepointer-5}{}\fi
\noindent\textit{Source reference.} Original reference stacks-properties-section-properties-morphisms: sources/stacks-algebraic-stacks-morphisms/pinned-original-source.tex, source lines 248--248. Navigation pointer only: the counted text refers to the Stacks item stacks-properties-section-properties-morphisms, which belongs to another chapter of the project and is not part of this card. The pointer carries the target's own label as its number because this chapter's staged snapshot does not contain the other chapter's tag table; its content is not reproduced here.\par\medskip
\begin{lemma}
\label{lemma-properties-diagonal-from-presentation}
Let $(U, R, s, t, c)$ be a groupoid in algebraic spaces such that
$s, t : R \to U$ are flat and locally of finite presentation.
Consider the algebraic stack $\mathcal{X} = [U/R]$ (see above).
\begin{enumerate}
\item If $R \to U \times U$ is separated, then
$\Delta_\mathcal{X}$ is separated.
\item If $U$, $R$ are separated, then $\Delta_\mathcal{X}$ is separated.
\item If $R \to U \times U$ is locally quasi-finite, then $\mathcal{X}$
is quasi-DM.
\item If $s, t : R \to U$ are locally quasi-finite, then
$\mathcal{X}$ is quasi-DM.
\item If $R \to U \times U$ is proper, then $\mathcal{X}$ is separated.
\item If $s, t : R \to U$ are proper and $U$ is separated, then
$\mathcal{X}$ is separated.
\item Add more here.
\end{enumerate}
\end{lemma}

\begin{proof}
Observe that the morphism $U \to \mathcal{X}$ is surjective, flat, and
locally of finite presentation by
Criteria for Representability, Lemma
\ref{criteria-lemma-flat-quotient-flat-presentation}.
Hence the same is true for $U \times U \to \mathcal{X} \times \mathcal{X}$.
We have the cartesian diagram
$$
\xymatrix{
R  = U \times_\mathcal{X} U \ar[r] \ar[d] & U \times U \ar[d] \\
\mathcal{X} \ar[r] & \mathcal{X} \times \mathcal{X}
}
$$
(see Groupoids in Spaces, Lemma
\ref{spaces-groupoids-lemma-quotient-stack-2-cartesian}).
Thus we see that $\Delta_\mathcal{X}$ has one of the properties listed in
Properties of Stacks, Section
\ref{stacks-properties-section-properties-morphisms}
if and only if the morphism $R \to U \times U$ does, see
Properties of Stacks, Lemma
\ref{stacks-properties-lemma-check-property-covering}.
This explains why (1), (3), and (5) are true.
The condition in (2) implies $R \to U \times U$ is separated
hence (2) follows from (1).
The condition in (4) implies the condition in (3)
hence (4) follows from (3).
The condition in (6) implies the condition in (5)
by Morphisms of Spaces, Lemma
\ref{spaces-morphisms-lemma-universally-closed-permanence}
hence (6) follows from (5).
\end{proof}

\end{document}
